% ============================================================
% style/commands.tex
%
% Удобные команды документа.
%
% Здесь находятся только пользовательские команды.
%
% НЕ здесь оформляются:
%   - заголовки \part, \chapter, \section, \subsection;
%   - внешний вид рамок rbox;
%   - оформление рисунков.
%
% Для этого используются отдельные файлы:
%   style/headings.tex
%   style/boxes.tex
%   style/figures.tex
%
% Зависимости:
%   \usepackage{xparse}
%   \usepackage{graphicx}
%   \usepackage{etoolbox}
%   \usepackage{ulem}
%   \usepackage{xcolor}
%   \usepackage{hyperref}
%
% Также предполагается, что уже определены:
%   - среда rbox из style/boxes.tex;
%   - цвета rbg, rc1--rc20 из style/colors.tex.
% ============================================================


% ------------------------------------------------------------
% 1. Большой подчёркнутый заголовок без номера
% ------------------------------------------------------------
%
% Использование:
%
%   \rtitle{Ознакомительная страница}
%
% По сути это:
%   \section*{...}

\NewDocumentCommand{\rtitle}{m}{%
    \phantomsection
    \section*{#1}
    \addcontentsline{toc}{section}{#1}
}

% ------------------------------------------------------------
% 2. Part-like страница
% ------------------------------------------------------------
%
% Использование:
%
%   \rpart{Семестр I}{Первый контакт}{}
%   \rpart{Семестр I}{Первый контакт}{image/pic.png}
%   \rpart[0.55\textwidth]{Семестр I}{Первый контакт}{image/pic.png}
%
% Аргументы:
%   [1] ширина картинки;
%   {2} главный текст;
%   {3} подпись под главным текстом;
%   {4} путь к картинке, можно оставить пустым.

\NewDocumentCommand{\rpart}{O{0.68\textwidth} m m m}{%
    \cleardoublepage

    \refstepcounter{part}%
    \setcounter{chapter}{0}%
    \setcounter{section}{0}%

    \phantomsection

    \ifstrempty{#3}{%
        \addcontentsline{toc}{part}{#2}%
    }{%
        \addcontentsline{toc}{part}{#2. #3}%
    }%

    \thispagestyle{empty}%
    \markboth{}{}%

    \null
    \vspace*{\fill}

    \begin{center}
        {\fontsize{36}{42}\selectfont\bfseries #2\par}

        \ifstrempty{#3}{}{%
            \vspace{0.8em}
            {\fontsize{17}{22}\selectfont\itshape #3\par}%
        }%

        \ifstrempty{#4}{}{%
            \vspace{0.5em}

            \includegraphics[
                width=#1,
                keepaspectratio
            ]{#4}%
        }%
    \end{center}

    \vspace*{\fill}
    \null

    \clearpage
}

% ------------------------------------------------------------
% 3. Короткая команда для семестра
% ------------------------------------------------------------
%
% Использование:
%
%   \rsem{I}{Первый контакт}{image/pic.png}
%
% Без картинки:
%
%   \rsem{I}{Первый контакт}{}
%
% С другой шириной картинки:
%
%   \rsem[0.55\textwidth]{I}{Первый контакт}{image/pic.png}
%
% По сути это сокращение для:
%
%   \rpart{Семестр I}{Первый контакт}{image/pic.png}

\NewDocumentCommand{\rsem}{O{0.68\textwidth} m m m}{%
    \rpart[#1]{Семестр #2}{#3}{#4}%
}


% ------------------------------------------------------------
% 4. Текстовый термин
% ------------------------------------------------------------
%
% Использование:
%
%   \rterm{Standing Waves}
%   Дальше идёт пояснение.
%
% Результат:
%
%   Standing Waves - дальше идёт пояснение.

\NewDocumentCommand{\rterm}{m}{%
    \par
    \noindent\textbf{#1} - %
}


% ------------------------------------------------------------
% 5. Центрированная строка
% ------------------------------------------------------------
%
% Использование:
%
%   \rcenter{Текст по центру}

\NewDocumentCommand{\rcenter}{m}{%
    \begin{center}
        #1
    \end{center}%
}


% ------------------------------------------------------------
% 6. Формульная вставка через rbox
% ------------------------------------------------------------
%
% Использование:
%
%   \begin{rformula}
%       \[
%           E = mc^2
%       \]
%   \end{rformula}
%
% С другим типом рамки:
%
%   \begin{rformula}{[]}
%       \[
%           E = mc^2
%       \]
%   \end{rformula}
%
% Типы рамок:
%   {[}   без правой стороны;
%   {[]}  полный прямоугольник;
%   {^}   без нижней стороны;
%   {v}   без верхней стороны;
%   {<}   левый угловой контур;
%   {>}   правый угловой контур.

\NewDocumentEnvironment{rformula}{O{<}}{%
    \begin{rbox}{#1}
    \centering
}{%
    \end{rbox}
}


% ============================================================
% 7. Быстрое цветовое выделение
% ============================================================
%
% Использование:
%
%   \rtxt{1}{красный текст}
%   \rhl{5}{текст с подложкой}
%   \rtag{12}{MVT}
%   \ruline{3}{подчёркнутый текст}
%
% Цвета:
%   1--10  красные оттенки;
%   11--20 сине-голубые / фиолетовые оттенки.
%
% Цветовые алиасы rc1--rc20 задаются в:
%   style/colors.tex
% ============================================================


% ------------------------------------------------------------
% 7.1. Внутренняя команда: применить цвет по номеру
% ------------------------------------------------------------
%
% Аргументы:
%   {1} номер цвета;
%   {2} команда-обработчик вида \handler{color}{text};
%   {3} текст.
%
% Если номер равен 0 или выходит за пределы 1--20,
% используется белый цвет.

\NewDocumentCommand{\rapplycolor}{m m m}{%
    \ifcase#1
        #2{white}{#3}%
    \or
        #2{rc1}{#3}%
    \or
        #2{rc2}{#3}%
    \or
        #2{rc3}{#3}%
    \or
        #2{rc4}{#3}%
    \or
        #2{rc5}{#3}%
    \or
        #2{rc6}{#3}%
    \or
        #2{rc7}{#3}%
    \or
        #2{rc8}{#3}%
    \or
        #2{rc9}{#3}%
    \or
        #2{rc10}{#3}%
    \or
        #2{rc11}{#3}%
    \or
        #2{rc12}{#3}%
    \or
        #2{rc13}{#3}%
    \or
        #2{rc14}{#3}%
    \or
        #2{rc15}{#3}%
    \or
        #2{rc16}{#3}%
    \or
        #2{rc17}{#3}%
    \or
        #2{rc18}{#3}%
    \or
        #2{rc19}{#3}%
    \or
        #2{rc20}{#3}%
    \else
        #2{white}{#3}%
    \fi
}


% ------------------------------------------------------------
% 7.2. Цветной текст
% ------------------------------------------------------------
%
% Использование:
%
%   \rtxt{5}{важный текст}

\NewDocumentCommand{\rtxtwith}{m m}{%
    {\textcolor{#1}{#2}}%
}

\NewDocumentCommand{\rtxt}{m m}{%
    \rapplycolor{#1}{\rtxtwith}{#2}%
}


% ------------------------------------------------------------
% 7.3. Цветное выделение фоном
% ------------------------------------------------------------
%
% Использование:
%
%   \rhl{5}{важный кусок}
%
% С другим внутренним отступом:
%
%   \rhl[3pt]{5}{важный кусок}
%
% Важно:
%   \rhl — inline-команда.
%   Она не предназначена для длинных абзацев.

\NewDocumentCommand{\rhlwith}{m m}{%
    {%
        \setlength{\fboxsep}{\rhlpad}%
        \colorbox{#1}{\textcolor{rbg}{#2}}%
    }%
}

\NewDocumentCommand{\rhl}{O{1.5pt} m m}{%
    \begingroup
        \def\rhlpad{#1}%
        \rapplycolor{#2}{\rhlwith}{#3}%
    \endgroup
}


% ------------------------------------------------------------
% 7.4. Маленькая цветная метка
% ------------------------------------------------------------
%
% Использование:
%
%   \rtag{12}{MVT}

\NewDocumentCommand{\rtagwith}{m m}{%
    {%
        \setlength{\fboxsep}{2pt}%
        \colorbox{#1}{%
            \textcolor{rbg}{\scriptsize\bfseries #2}%
        }%
    }%
}

\NewDocumentCommand{\rtag}{m m}{%
    \rapplycolor{#1}{\rtagwith}{#2}%
}


% ------------------------------------------------------------
% 7.5. Цветное подчёркивание
% ------------------------------------------------------------
%
% Использование:
%
%   \ruline{3}{подчёркнутый текст}
%
% Команда использует механизм пакета ulem:
%   \markoverwith
%   \ULon

\NewDocumentCommand{\rulinewith}{m m}{%
    {%
        \bgroup
            \markoverwith{%
                \textcolor{#1}{\rule[-0.55ex]{2pt}{0.45pt}}%
            }%
            \ULon{#2}%
        \egroup
    }%
}

\NewDocumentCommand{\ruline}{m m}{%
    \rapplycolor{#1}{\rulinewith}{#2}%
}

% ------------------------------------------------------------
% 8. Вспомогательные команды
% ------------------------------------------------------------

\newcommand{\practice}[3]{%
    \phantomsection
    \addcontentsline{toc}{chapter}{Практическое занятие №#1. #3}%
    \begin{center}
        {\Large\bfseries Практическое занятие №#1}\\[0.3em]
        {\small #2}\\[0.6em]
        {\large\bfseries #3}
    \end{center}
    \vspace{0.4em}
    \hrule
    \vspace{1.2em}
}